\section{Membiasakan Berpirik Kritis dan Mencintai IPTEK}

\subsection{Dalil Berpikir Kritis dalam Al-Quran}

\begin{itemize}
\item \textit{Inna fī khalqis-samāwāti wal-arḍi wakhtilāfil-laili wan-nahāri la'āyātil li'ulil-albāb.} \\
  \textbf{Artinya: }Sesungguhnya dalam penciptaan langit dan bumi dan pergantian malam dan siang terdapat tanda-tanda (kebesaran Allah) bagi orang yang berakal. \textbf{(Al-Imran (190))}
  
\item \textit{Allażīna yażkurūnallāha qiyāmaw wa qu‘ūdaw wa ‘alā junūbihim wa yatafakkarūna fī khalqis-samāwāti wal-arḍ, rabbanā mā khalaqta hāżā bāṭilā, subḥānaka fa qinā ‘ażāban-nār.} \\
  \textbf{Artinya: }(Yaitu) orang-orang yang mengingat Allah sambil berdiri, duduk, atau dalam keadaan berbaring, dan mereka memikirkan tetang penciptaan langit dan bumi (seraya berkata) ``Ya tuhan kami, tidaklah engkau menciptakan semua ini sia-sia''. Mahasuci engkau lindungilah kami dari azab neraka. \textbf{Al-Imran (191)}

\item \textit{Yā ma‘syaral-jinni wal-insi inistaṭa‘tum an tanfużū min aqṭāris-samāwāti wal-arḍi fanfużū, lā tanfużūna illā bisulṭān.} \\
  \textbf{Artinya: }Wahai golongan jin dan manusia! jika kamu sanggup menembus (melintasi) penjuru langit dan bumi, maka tembuslah, kamu tidak akan mampu menembusnya kecuali dengan kekuatan (dari Allah). \textbf{Ar-Rahman (33)}
\end{itemize}

\subsection{Dalil Berpikir Kritis dalam Hadits}
\begin{itemize}
\item \textit{‘An Abī Hurairata raḍiyallāhu ‘anhu qāla: qāla Rasūlullāhi ṣallallāhu ‘alaihi wa sallam: Kafā bil-mar'i każiban ay-yuḥaddisa bikulli mā sami‘a.} \\
  \textbf{Artinya: }Dari Abu Hurairah R.A ia berkata: Rasulullah bersabda: ``Cukuplah seseorang disebut pendusta yaitu orang yag mengatakan (membicarakan) semua yang ia dengar.'' \textbf{(HR Muslim)}.

\item \textit{‘An Abī Żarrin raḍiyallāhu ‘anhu qāla: qāla Rasūlullāhi ṣallallāhu ‘alaihi wa sallam: Tafakkarū fī khalqillāhi wa lā tafakkarū fī żātillāhi fata hlikū.} \\
  \textbf{Artinya: }Dari Abi Dzar R.A ia berkata: Nabi SAW bersabda: ``Pikirkalah mengenai segala sesuatu ciptaan Allah. Tetapi janganlah kalian memikirkan tetang zat Allah, karena kalian akan rusak.'' \textbf{(HR Abu Syaih)}
\end{itemize}
